\documentclass[10pt,twocolumn,letterpaper]{article}

% ---------------------------------------------------------------
% Include basic CVPR package in rebuttal mode

\usepackage[rebuttal]{cvpr}

% ---------------------------------------------------------------
% Other packages

% Commonly used abbreviations (\eg, \ie, \etc, \cf, \etal, etc.)
\usepackage{eccvabbrv}

% Include other packages here, before hyperref.
\usepackage{graphicx}
\usepackage{booktabs}
\usepackage{xcolor}
\usepackage{amssymb}
\usepackage{amsmath}

% The "axessiblity" package can be found at: https://ctan.org/pkg/axessibility?lang=en
\usepackage[accsupp]{axessibility}  % Improves PDF readability for those with disabilities.


% ---------------------------------------------------------------
% Hyperref package
% If you comment hyperref and then uncomment it, you should delete
% egpaper.aux before re-running latex.  (Or just hit 'q' on the first latex
% run, let it finish, and you should be clear).

\definecolor{eccvblue}{rgb}{0.12,0.49,0.85}
\definecolor{revred}{RGB}{180,35,35}
\definecolor{revblue}{RGB}{35,95,180}
\definecolor{revgreen}{RGB}{0,120,80}
\newcommand{\RQ}{\textcolor{revred}{\textbf{Q4FJ}}}
\newcommand{\RN}{\textcolor{revblue}{\textbf{Nj5w}}}
\newcommand{\RP}{\textcolor{revgreen}{\textbf{PPhm}}}
\newcommand{\good}[1]{\textcolor{revgreen}{\textbf{#1}}}
\newcommand{\crit}[1]{\textcolor{revred}{\textbf{#1}}}
\newcommand{\src}[1]{\textcolor{revblue}{\footnotesize\,(#1)}}
\newcommand{\tightpara}[1]{\vspace{-0.9em}\paragraph{#1}\vspace{-0.25em}}
\setlength{\parskip}{0pt}
\setlength{\parindent}{0.55em}
\usepackage[pagebackref,breaklinks,colorlinks,citecolor=eccvblue]{hyperref}


% ---------------------------------------------------------------
% If you wish to avoid re-using figure, table, and equation numbers from
% the main paper, please uncomment the following and change the numbers
% appropriately.
%\setcounter{figure}{2}
%\setcounter{table}{1}
%\setcounter{equation}{2}

% If you wish to avoid re-using reference numbers from the main paper,
% please uncomment the following and change the counter for `enumiv' to
% the number of references you have in the main paper (here, 6).
%\let\oldthebibliography=\thebibliography
%\let\oldendthebibliography=\endthebibliography
%\renewenvironment{thebibliography}[1]{%
%     \oldthebibliography{#1}%
%     \setcounter{enumiv}{6}%
%}{\oldendthebibliography}


% ---------------------------------------------------------------
% TODO REBUTTAL: Please enter paper ID
\def\paperID{4870} % *** Enter the paper ID here
\def\confName{ECCV}
\def\confYear{2026}

\begin{document}
\setlength{\linenumbersep}{0.35cm}
\emergencystretch=1.2em
\setlength{\abovedisplayskip}{2pt}
\setlength{\belowdisplayskip}{2pt}
\setlength{\abovedisplayshortskip}{1pt}
\setlength{\belowdisplayshortskip}{1pt}

% ---------------------------------------------------------------
% TODO REBUTTAL: Paper title is not required
%\title{\LaTeX\ Guidelines for Author Response}  % **** Enter the paper title here
\title{Author Response}

\maketitle
\thispagestyle{empty}
\appendix
\vspace{-2.2em}


% ---------------------------------------------------------------
% TODO REBUTTAL: Enter your response below
\noindent\textbf{We thank all reviewers for the constructive feedback.}
\RQ{} recognized SRUM as an elegant self-evaluation method with strong gains; \RN{} highlighted the real UMM understanding-generation gap, global-local rewards, Fig.~3/4 analyses, and practical overhead; \RP{} appreciated the timely self-rewarding paradigm and compelling gains/generalization. We will add complete loss/parameter scope, self-scoring protocol, reward-map pseudocode, prompt templates, and saved-candidate diagnostics.

\tightpara{\RQ: objective and training setting.}
There is no hidden objective beyond Eq.~(3). ``KL constraint'' in Sec.~4.3 was inaccurate: it means reference constraint \(L_{\rm ref}\), not an extra KL term. The full loss is \(L=L_r+\lambda_cL_{\rm ref}\), \good{\(\lambda_c=0.5\)}, where \(L_r=E[\alpha R\odot\|v_\theta-(\epsilon-x^{gt}_0)\|_2^2]\) and \(L_{\rm ref}=E[\|v_\theta-(\epsilon-x^{gt}_0)\|_2^2]\). SRUM is \emph{offline}: generate \good{6K} candidates, freeze the UMM understanding side as evaluator (\emph{no evaluator gradients}), cache rewards, then update Bagel stage-4 trainable generation/flow parameters. This separates reward estimation from optimization and avoids evaluator drift. We will state this scope, recipe (AdamW, lr \good{\(2.5{\times}10^{-5}\)}, warmup 500, grad clip 1.0), and sampling defaults \src{Main Sec.~3.3/Fig.~2; Supp. A}.

\tightpara{\RQ: self-evaluation details and validation.}
Appendix prompts specify the scoring interface. Region proposals become IDs/boxes; the frozen UMM identifies each object and assigns a local \([-1,1]\) relevance/quality score with a reason. A global-layout prompt scores whole-image alignment and is normalized to \([0,1]\) as \(\alpha\). Local scores fill \(R\), and \(\alpha\) gates \(R\). Rewards come from the UMM's own judgments, not human labels or external reward models; prompts standardize parsing only and were not tuned on downstream validation \src{Main Sec.~3.1-3.2; Supp. Tab.~8-10}. We will add diagnostics from existing candidates; Fig.~3/4 show causal utility by component, not only final metrics.

\tightpara{\RQ/\RN: external evaluators and calibration.}
We agree Table~3 deserves more analysis. \RQ asked about external models and a frozen UMM evaluator; our setting is exactly a frozen internal UMM scorer during reward generation. There is no hidden loss and no online evaluator update: ``KL'' will be corrected to \(L_{\rm ref}\), and reward scoring uses \emph{no evaluator gradients}. Calibration is computed by rescoring the same \good{1K} saved image/prompt pairs with internal UMM and Qwen-32B under the same JSON rubric, aggregating region rewards to image-level rewards, then computing Spearman \(\rho\) with a held-out final-evaluator ranking. Higher \(\rho\) means the reward better ranks candidates by final generation quality, testing reward usefulness rather than standalone VLM accuracy. This explains why the internal scorer can be the better training signal even if Qwen-32B is a stronger generic VLM \src{Main Tab.~3}.
\vspace{-0.45em}
\begin{center}
\scriptsize
\setlength{\tabcolsep}{2pt}
\renewcommand{\arraystretch}{0.9}
\begin{tabular*}{0.94\columnwidth}{@{\extracolsep{\fill}}lcc@{}}
\toprule
\textbf{Diagnostic} & \textbf{Ours} & \textbf{Ref.}\\
\midrule
Region validity & \good{92.3\%} & --\\
Accept/reject reward & \good{0.67} & \crit{-0.31}\\
Calib. \(\rho\) & \good{.58} & \crit{.43}\\
Reward-source score & \good{88.37} & \crit{88.01}\\
\bottomrule
\end{tabular*}
\end{center}
\vspace{-0.8em}

\tightpara{\RQ/\RN: ablation consistency and efficiency.}
The gains are not from candidate data alone: SFT on the same self-generated data is weaker than SRUM, so the reward weighting is the critical factor. Main T2I-CompBench improves \good{\(84.46\!\to\!88.37\)} (\good{+3.91}); OOD T2I-ReasonBench improves \good{\(43.82\!\to\!46.75\)} (\good{+2.93}); traditional metrics also improve (\good{+2.52/+2.18} avg.). We will fix the text/table inconsistency: ``w/o local reward'' belongs to Fig.~3(left), not Table~3. These gains use one offline round, not recursive self-training. The cost is small: scoring \good{6K} candidates costs \good{\(<4\)} GPU-hours, and training is nearly SFT-cost (\good{12.5} vs. 12.4 GPU-hours), without manual cleaning \src{Main Fig.~3/Tab.~1-3/6}.

\tightpara{\RN/\RP: SAM dependency and self-rewarding scope.}
We agree that ``self-contained'' was overstated for Bagel. SAM proposes spatial supports for Bagel, while BLIP3o can rely on native grounding; in both cases the UMM verifies boxes, matches prompt-regions, and assigns all semantic rewards. We will revise ``no external model'' to ``no external reward model or human labels.'' Replacing SAM with weaker UMM-only grounding retains most gains in our diagnostic (\(\sim\)\good{87.9} overall), suggesting SAM is a localization aid rather than the semantic reward source. SRUM is also not only a bbox method: global reward covers layout, style, and knowledge consistency, while local reward refines object fidelity \src{Main Sec.~3.1-3.2; Fig.~4; Supp. A}.

\tightpara{\RP/\RN: evaluation bias, degeneration, and generality.}
QwenVL is only an evaluation metric and is never used in SRUM training. Table~2's traditional metrics, GenEval, WISE, and T2I-ReasonBench reduce single-evaluator bias concerns. We did not observe collapse in one-round offline post-training: understanding is preserved (MMBench \good{\(85.0/84.8\)}, MM-Vet \good{\(67.2/67.0\)}) and MMVP improves \good{\(69.3\!\to\!70.0\)}. Since we evaluate two 7B UMMs, Bagel and BLIP3o, we will tone down broader generality claims beyond tested families/scales \src{Main Sec.~4.1; Tab.~4-6}.

\tightpara{Writing corrections.}
We will fix the missing reference, typo near line 351, and Fig.~3/Table~3 inconsistency. We will sharpen concurrent-work discussion: unlike rule-level or RL-only self-improvement, SRUM builds dense global-local reward maps and directly weights flow training. We will also highlight Appendix D limitations: wrong language reasoning, causal multi-image generation, and aesthetics-focused prompts \src{Supp. D}.

\end{document}
